% Gesamttest «Potenz- und Wurzelfunktionen» — Quelle des PDFs (python3 scripts/build-lp-pdf.py).
% Musterlösung und Raster: bewertungspaket.tex. Alle Werte mit python3 nachgerechnet (03.10.2026).
% Kein \lt/\gt — pdfLaTeX kennt sie nicht, < und > tun es.
\documentclass[11pt]{article}
\usepackage{../lp-druck}
\renewcommand{\lpname}{Potenz- und Wurzelfunktionen}
\renewcommand{\lpteilgebiet}{3.2}
\renewcommand{\lpbereich}{Schwerpunktfach}
\renewcommand{\lpfuss}{Gesamttest · 26 Punkte}
\begin{document}
\lpkopf{Gesamttest}{%
  \vspace{4pt}\noindent\begin{tabularx}{\linewidth}{@{}X X@{}}
  Code (kein Name): \hrulefill & Punkte: \hrulefill\ / 26
  \end{tabularx}}

\begin{lpinfo}{Anleitung}
\textbf{Zeit:} rund 30 Minuten. \textbf{Hilfsmittel:} kein Taschenrechner — die Kompetenz des Teilgebiets 3.2 trägt im Lehrplan den Vermerk «auch ohne Hilfsmittel», und die Aufgaben aus Kapitel 3 (Teilgebiet 3.1) sind ebenfalls ohne Rechner lösbar. Lineal und Bleistift sind erlaubt.
Schreib zu jeder Aufgabe den \textbf{Rechenweg} auf — bewertet wird der Weg.
Danach: Lösung scannen oder fotografieren (ohne Namen und Standort) und mit dem \emph{Bewertungspaket} bewerten lassen.
\end{lpinfo}

\lpteil{Teil A · Potenzfunktionen}{Kapitel 1 und 2 · 9 Punkte}

\begin{lpaufgabe}{G1}{Werte und Symmetrie}{3 P}
Gegeben ist $f(x) = 2x^{4}$.
\begin{enumerate}[label=(\alph*),leftmargin=*,itemsep=1pt]
\item Berechne $f(-2)$ und $f(0.5)$.
\item Welche Symmetrie hat der Graph? Weise sie mit $f(-x)$ nach.
\end{enumerate}
\lpplatz{6}
\end{lpaufgabe}

\begin{lpaufgabe}{G2}{Von der Kurve zur Gleichung}{3 P}
\begin{minipage}[t]{0.52\linewidth}
Gib die Gleichung der abgebildeten Potenzfunktion $f(x) = a \cdot x^{n}$ an.
Die markierten Punkte liegen auf Gitterpunkten. Begründe, woran du $n$ erkennst.
\lpplatz{6}
\end{minipage}\hfill
\begin{minipage}[t]{0.42\linewidth}\vspace{0pt}\centering
\begin{tikzpicture}
\begin{axis}[width=6.4cm,height=7.6cm,axis lines=middle,xmin=-2.4,xmax=2.4,ymin=-9,ymax=9,
  xtick={-2,-1,1,2},ytick={-8,-6,-4,-2,2,4,6,8},grid=both,grid style={lplinie},tick label style={font=\scriptsize},
  xlabel=$x$,ylabel=$y$,
  yticklabel style={fill=white,inner sep=1pt,xshift=-2pt},xticklabel style={fill=white,inner sep=1pt}]
\addplot[lpblau,very thick,domain=-2.05:2.05,samples=120] {-x^3};
\addplot[only marks,mark=*,color=lpgruen] coordinates {(1,-1) (2,-8)};
\end{axis}
\end{tikzpicture}
\end{minipage}
\end{lpaufgabe}

\begin{lpaufgabe}{G3}{Eine Hyperbel}{3 P}
Gegeben ist $g(x) = -\dfrac{2}{x^{2}}$.
\begin{enumerate}[label=(\alph*),leftmargin=*,itemsep=1pt]
\item Gib Definitionsmenge und die beiden Asymptoten an.
\item Wo liegen die beiden Äste? Begründe mit der Parität des Exponenten und dem Vorzeichen von $a$.
\end{enumerate}
\lpplatz{6}
\end{lpaufgabe}

\lpteil{Teil B · Verschieben und umkehren}{Kapitel 3 und 4 · 11 Punkte}

\begin{lpaufgabe}{G4}{Verschoben: Nullstellen und Asymptoten}{4 P}
\begin{enumerate}[label=(\alph*),leftmargin=*,itemsep=1pt]
\item Berechne die Nullstellen von $h(x) = (x-1)^{4} - 16$ und begründe in einem Satz, warum es genau so viele sind.
\item Gib für $g(x) = \dfrac{2}{x-3} + 1$ die Definitionsmenge und beide Asymptoten an.
\end{enumerate}
\lpplatz{7}
\end{lpaufgabe}

\begin{lpaufgabe}{G5}{Umkehren}{4 P}
\begin{minipage}[t]{0.52\linewidth}
Gegeben ist $f(x) = x^{3} - 2$.
\begin{enumerate}[label=(\alph*),leftmargin=*,itemsep=1pt]
\item Bestimme $f^{-1}$ mit Rechenweg.
\item Auf dem Graphen von $f$ liegt $P(2 \mid 6)$. Welcher Punkt liegt auf $f^{-1}$?
\item Skizziere beide Kurven und die Spiegelachse ins Koordinatensystem.
\end{enumerate}
\end{minipage}\hfill
\begin{minipage}[t]{0.42\linewidth}\vspace{0pt}\centering
\begin{tikzpicture}
\begin{axis}[width=6.4cm,height=6.4cm,axis lines=middle,xmin=-7,xmax=7,ymin=-7,ymax=7,
  xtick={-6,-4,-2,2,4,6},ytick={-6,-4,-2,2,4,6},grid=both,grid style={lplinie},tick label style={font=\scriptsize},
  xlabel=$x$,ylabel=$y$,axis equal image,
  yticklabel style={fill=white,inner sep=1pt,xshift=-2pt},xticklabel style={fill=white,inner sep=1pt}]
\end{axis}
\end{tikzpicture}
\end{minipage}
\lpplatz{6}
\end{lpaufgabe}

\begin{lpaufgabe}{G6}{Warum einschränken?}{3 P}
Warum muss man $y = x^{4}$ einschränken, bevor man sie umkehrt, $y = x^{5}$ aber nicht?
Belege beides mit je einem Zahlenbeispiel.
\lpplatz{6}
\end{lpaufgabe}

\lpteil{Teil C · Wurzelfunktionen}{Kapitel 5 · 6 Punkte}

\begin{lpaufgabe}{G7}{Zwei Wurzelkurven}{4 P}
\begin{enumerate}[label=(\alph*),leftmargin=*,itemsep=1pt]
\item Gegeben ist $k(x) = 2\sqrt{x+3} - 4$. Gib Definitionsmenge und Startpunkt an und berechne die Nullstelle.
\item Und $m(x) = \sqrt[3]{x+1} - 2$? Gib die Definitionsmenge an und sag in einem Satz, warum sie anders aussieht als bei $k$.
\end{enumerate}
\lpplatz{7}
\end{lpaufgabe}

\begin{lpaufgabe}{G8}{Der Grösse nach}{2 P}
Ordne $\sqrt{x}$, $x$ und $x^{2}$ der Grösse nach für $0 < x < 1$ und belege die Reihenfolge mit einem Zahlenbeispiel.
\lpplatz{5}
\end{lpaufgabe}
\end{document}
